\documentclass[11pt,a4paper]{article}
\usepackage[utf8]{inputenc}
\usepackage[T1]{fontenc}
\usepackage{amsmath,amssymb}
\usepackage[margin=2.4cm]{geometry}
\usepackage{booktabs}
\usepackage{hyperref}
\hypersetup{colorlinks=true,linkcolor=blue,citecolor=blue,urlcolor=blue}
\usepackage{natbib}

\title{\bf The matter dilution exponent is not identifiable\\
from compressed CMB priors:\\
a referent ambiguity, not a precision limit}
\author{\'Edouard Lantenois\thanks{\texttt{github.com/Dantenos}}}
\date{August 2026}

\begin{document}
\maketitle

\begin{abstract}
\noindent
Consider the one-parameter family in which dark matter dilutes as
$\rho_{\rm dm}\propto a^{-3+\epsilon}$ with a genuine cosmological constant. The Bianchi
identity forces this background to be that of a fluid with constant equation of state
$w=-\epsilon/3$, so the family is not an alternative to $\Lambda w$DM but a rewriting of it.
Compressed CMB distance priors $(\mathcal{R}, \ell_A, \omega_b)$, together with the standard
fitting formulas for $r_d$, $r_\ast$ and $z_\ast$, all require a \emph{single} value of
$\omega_m$. When $\epsilon\neq0$ there are two --- the present-day label and the value at
recombination --- and nothing in the data selects one. We show on identical data
(Pantheon+, DESI DR2 BAO, Planck 2018 distance priors, SH0ES) that this bookkeeping choice
alone moves the recovered exponent from $\epsilon=+0.0070$ to $\epsilon=-0.0120$, a range of
$0.019$, while every configuration quotes $\sigma\simeq0.002$. We then show that this is
\emph{not} an information-loss effect: the compressed priors constrain $\omega_c$ to
$\sigma=0.0010$, the same value the full Planck TT,TE,EE$+$low-$\ell$ likelihood delivers
under identical fixed parameters, a ratio of $1.00$ (bounded to $]0.5,1.25[$ by grid
quantisation). The ambiguity is therefore one of \emph{referent}, not of \emph{precision},
and the spread is nineteen times the available precision. Consistently, published
determinations of this exponent obtained with modified Boltzmann codes and full CMB
likelihoods cluster at $\epsilon = -0.00064\pm0.00045$ and span $0.005$, while
compressed-prior determinations span $0.016$. We report one further trap: applying
$\epsilon$ to \emph{total} matter rather than to dark matter alone creates an internal
inconsistency worth $7.5$--$15.1\sigma$ in $\ell_A$, because the same likelihood then treats
baryons as standard in $R_b$ and in the $\omega_b$ prior. We claim no detection. We claim
that a published $\epsilon$ obtained from compressed priors carries an unquoted systematic
of order $10^{-2}$ unless the choice of $\omega_m$ is stated.
\end{abstract}

\section{The family, and why it has only one background}
\label{sec:famille}

Write the dark sector as a matter component whose density obeys
\begin{equation}
\rho_{\rm dm}(a) = \rho_{{\rm dm},0}\,a^{-3+\epsilon},
\qquad
\rho_\Lambda = \text{const},
\qquad
\rho_b \propto a^{-3},\quad \rho_r\propto a^{-4}.
\label{eq:famille}
\end{equation}
It is tempting to read \eqref{eq:famille} as \emph{pressureless matter that is not
conserved}, the deficit being absorbed by nothing in particular. That reading is not
available. With $\Lambda$ genuinely constant, $\nabla_\mu T^{\mu\nu}=0$ applied to the total
stress tensor forces the matter sector to carry a pressure
\begin{equation}
p_{\rm dm} = -\tfrac{\epsilon}{3}\,\rho_{\rm dm},
\qquad\text{i.e.}\qquad
w_{\rm dm} = -\epsilon/3 .
\label{eq:bianchi}
\end{equation}
``Non-conserved pressureless matter with constant $\Lambda$'' and ``conserved matter with
constant $w=-\epsilon/3$ and constant $\Lambda$'' are the same model written twice. At
background level there is one model, not two.

This is not new. \citet{lima1996} give, for adiabatic matter creation with rate parameter
$\beta$, the effective adiabatic index $\gamma_\ast=\gamma(1-\beta)$ and the scaling
$\rho=\rho_0(a_0/a)^{3\gamma(1-\beta)}$, with creation pressure $p_c=-\beta\gamma\rho$. For
dust ($\gamma=1$) this is exactly \eqref{eq:famille}--\eqref{eq:bianchi} with
$\epsilon=3\beta$. Two cautions, which we state because they are easy to lose: the identity
$\gamma_\ast=\gamma(1-\beta)$ holds only for \emph{constant} $\beta$, and the map
$\epsilon=3\beta$ is \emph{species-dependent} --- for radiation ($\gamma=4/3$) the same
$\beta$ gives $\rho_r\propto a^{-4(1-\beta)}$. The map may not be carried across the
radiation era.

We also note that the symbol $\epsilon$ is used in a nearby but different literature ---
that of \emph{decaying vacuum}, e.g.\ \citet{alcaniz2005} --- where $\Lambda$ is not
constant. The family \eqref{eq:famille} is the constant-$\Lambda$ one, and at background
level it coincides with the $\Lambda w$DM model constrained by \citet{kumar2025}.

\section{Where the ambiguity enters}
\label{sec:ou}

Compressed analyses of this family use three ingredients, and every one of them takes
$\omega_m$ as an input:
\begin{enumerate}
\item the shift parameter $\mathcal{R}=\sqrt{\Omega_m H_0^2}\,D_A(z_\ast)/c$, whose
      $\sqrt{\Omega_m}$ prefactor is derived assuming $\rho_m\propto a^{-3}$;
\item the fitting formulas for the drag and last-scattering sound horizons $r_d$, $r_\ast$
      \citep{eh1998}, calibrated on $\Lambda$CDM grids;
\item the fitting formula for $z_\ast$ \citep{hs1996}.
\end{enumerate}
In $\Lambda$CDM the input is unambiguous. In \eqref{eq:famille} it is not: the label
$\Omega_m$ and the density that actually governs pre-recombination physics differ by
\begin{equation}
\frac{\Omega_m^{\rm rec}}{\Omega_m^{\rm label}} = a_\ast^{\,\epsilon},
\qquad a_\ast\simeq 1/1091,
\qquad \ln a_\ast \simeq -6.995 .
\label{eq:ratio}
\end{equation}
For $|\epsilon|=0.005$ the two differ by $3.6\%$. An analyst must choose, and the literature
does not say which. \citet{elgaroy2007} warned that $\mathcal{R}$ ``is not directly
measurable from the cosmic microwave background, and its value is usually derived from the
data assuming a spatially flat cosmology with dark matter and a cosmological constant'';
we quantify what that warning costs in this family.

\section{The measurement: four bookkeeping choices, one dataset}
\label{sec:mesure}

We use Pantheon+ (1580 SNe, $z>0.01$, calibrators excluded), the 13 DESI DR2 BAO
measurements, Planck 2018 distance priors, and the SH0ES $H_0$ prior. Four configurations
differ \emph{only} in bookkeeping:

\begin{itemize}
\item[(a)] \textbf{label}: the fitting formulas and $\mathcal{R}$ are fed the present-day
  $\Omega_m$;
\item[(b)] \textbf{coherent}: they are fed $\Omega_m a_\ast^{\epsilon}$, the density at
  recombination;
\item[(c)] \textbf{direct, no $\mathcal{R}$}: $r_s$ is obtained by integrating
  $\int c_s\,\mathrm{d}a/(a^2H)$ in the model's own background, and $\mathcal{R}$ is dropped,
  leaving $(\ell_A,\omega_b)$ with the marginal $2\times2$ covariance;
\item[(d)] \textbf{direct, $\mathcal{R}$ coherent}: as (c) but retaining $\mathcal{R}$ fed
  the recombination density.
\end{itemize}

The direct integrator is validated against the CAMB table it replaces: with
$N_{\rm eff}=3.046$ and the massive-neutrino contribution removed from the matter budget,
the ratio $r_s^{\rm direct}/r_s^{\rm CAMB}$ varies by $1.3\times10^{-6}$ across the whole
$(\omega_b,\omega_m)$ grid, i.e.\ $0.01\sigma$ in $\ell_A$. Its inferred drag redshift,
recovered by inverting the integrator against CAMB's $r_{\rm drag}$, lands in
$[1057.8,1062.1]$ against Planck's $1059.9$ --- a check we did not impose.

\begin{table}[t]
\centering
\begin{tabular}{lrrr}
\toprule
configuration & $\epsilon$ & $\sigma$ (opposable) & $\Delta\chi^2$ vs $\Lambda$CDM \\
\midrule
(a) label                 & $+0.0070$ & $2.33$ & $+9.36$ \\
(b) coherent              & $-0.0040$ & $2.00$ & $+9.36$ \\
(c) direct, no $\mathcal{R}$ & $-0.0120$ & $3.14$ & $+9.87$ \\
(d) direct, $\mathcal{R}$ coherent & $-0.0040$ & $2.00$ & $+9.90$ \\
\midrule
\textbf{range} & \textbf{0.0190} & & \\
\bottomrule
\end{tabular}
\caption{Same data, same model, four bookkeeping choices. The quoted $\sigma$ is the smaller
of the profile half-width and the Wilks value $\sqrt{\Delta\chi^2}$; where the two disagree
the profile is non-parabolic and the local $\sigma$ is not opposable.}
\label{tab:quatre}
\end{table}

Table~\ref{tab:quatre} is the central measurement. The range is $0.019$; each configuration
quotes $\sigma\simeq0.002$. Two remarks on honesty. First, configuration (c) reports a local
significance of $6.0\sigma$ that we do \emph{not} quote: its profile is flat
($\Delta\chi^2(2\sigma_{\rm local})/4 = 0.39$ instead of $1$) and the likelihood-ratio value
is $3.14$. Second, a decomposition of (c)'s gain shows $126\%$ of it coming from the SH0ES
term alone, with the supernovae worsening by $1.72$ and the BAO by $0.96$: that configuration
buys $H_0$, it does not measure dilution. We report it because it belongs to the range, not
because we believe it.

\section{It is not information loss}
\label{sec:pouvoir}

The natural explanation --- compressed priors carry only geometry, the full likelihood also
carries the acoustic peak \emph{heights} that pin $\omega_m$ --- is wrong, and we tested it
rather than assert it.

The premise is correct: evaluating the compressed $\chi^2$ at $\ln(10^{10}A_s)=3.000$ and
$3.100$ returns $2.575399113$ both times, to machine precision. It carries no amplitude
information whatsoever.

The conclusion does not follow. Profiling $\omega_c$ with CAMB and the Planck
TT,TE,EE$+$low-$\ell$ likelihood at fixed $\theta_\ast$, $\omega_b$, $n_s$, $\tau$, and
profiling $\ln(10^{10}A_s)$ at each point, gives $\sigma(\omega_c)=0.0010$. The same
profile through the compressed priors gives $\sigma(\omega_c)=0.0010$. The ratio is
$1.00$; propagating the grid quantisation bounds it in $]0.5,1.25[$, far from any
information-loss signature. As a consistency check, our $\sigma$ is $1.20$ times smaller
than the marginalised Planck 2018 value $\omega_c=0.1200\pm0.0012$ \citep{planck2018}, which
is what fixing four parameters should do; and the full $\chi^2$ at its $\Lambda$CDM optimum
reproduces our frozen anchor to $0.005$.

So both determinations of $\omega_m$ are equally sharp --- $0.7\%$ --- and the spread across
bookkeeping choices is $0.019$, \emph{nineteen times} the available precision.

\begin{equation}
\boxed{\;\text{The ambiguity is of referent, not of precision.}\;}
\end{equation}

$\mathcal{R}$ and the fitting formulas demand one $\omega_m$; the family supplies two;
the data adjudicate neither. Four analysts with identical data and identical $0.7\%$
constraints obtain answers spanning $0.019$ --- not by noise, but by choice.

\section{The baryon trap}
\label{sec:baryons}

A second failure mode is worth isolating because it is easy to fall into and it is large.
If $\epsilon$ is applied to \emph{total} matter --- i.e.\ if $\Omega_m$ in \eqref{eq:famille}
is read as baryons plus dark matter --- then baryons dilute anomalously too. But the same
likelihood typically (i) compares $\omega_b$ to a prior measured under $\rho_b\propto a^{-3}$
and (ii) integrates the sound horizon with $R_b=(3\omega_b/4\omega_\gamma)\,a$, a form that
assumes the same. Consistency would require $R_b\propto a^{1+\epsilon}$. Measuring the
difference:

\begin{table}[h]
\centering
\begin{tabular}{lrr}
\toprule
 & $|\epsilon|=0.003$ & $|\epsilon|=0.006$ \\
\midrule
$r_s$ under $R_b\propto a$ vs $a^{1+\epsilon}$ & $7.5\sigma$ in $\ell_A$ & $15.1\sigma$ \\
$\omega_b$(recombination) vs the Planck prior & $4.8\sigma$ & $8.1\sigma$ \\
\bottomrule
\end{tabular}
\caption{Cost of applying $\epsilon$ to total matter while treating baryons as standard
elsewhere in the same likelihood.}
\label{tab:baryons}
\end{table}

Both exceed anything the analysis is trying to measure. The remedy is to restrict $\epsilon$
to the dark sector, as \citet{kumar2025} and \citet{yaoliu2025} do. We note, against our own
convenience, that the often-invoked argument ``BBN forbids anomalous baryon dilution'' is
weaker than it sounds: requiring the BBN and CMB determinations of $\omega_b$ to agree to
$\sim2.5\%$ across the $\simeq12.5$ e-folds between the two epochs bounds only
$|\epsilon_b|\lesssim2\times10^{-3}$, and we could find no published constraint on the baryon
dilution exponent itself. The defect above is one of \emph{internal} inconsistency, not a
violation of an external bound.

\section{Confrontation with the published determinations}
\label{sec:litterature}

We collected every determination of this exponent we could verify at primary source,
converting each into the convention of \eqref{eq:famille} by a declared rule, and split them
by method.

\begin{table}[t]
\centering
\begin{tabular}{llrr}
\toprule
determination & method & $\epsilon$ & $\sigma$ \\
\midrule
\multicolumn{4}{l}{\emph{full likelihood, modified Boltzmann code}}\\
\citet{tsiapi2019}, $\Lambda(H)$CDM2, Planck alone & CAMB, Planck 2015 & $+0.00148$ & $0.00252$\\
\citet{tsiapi2019}, $\Lambda(H)$CDM2, joint        & CAMB, Planck 2015 & $-0.00020$ & $0.00181$\\
\citet{tsiapi2019}, $\Lambda(H)$CDM1, joint$^\dagger$ & CAMB, Planck 2015 & $+0.00302$ & $0.00151$\\
\citet{kumar2025}, PL18$+$DESI DR2                 & CAMB, Planck 2018 & $-0.00194$ & $0.00096$\\
\citet{li2025}, Planck$+$DESI DR2$+$DESY5          & CAMB, NPIPE       & $-0.00126$ & $0.00075$\\
\citet{li2024}, CMB$+$DESI DR1$+$DESY5$^\ddagger$  & CosmoMC           & $-0.00025$ & $0.00092$\\
\midrule
\multicolumn{2}{l}{\textbf{weighted mean}} & $\mathbf{-0.00064}$ & $\mathbf{0.00045}$\\
\multicolumn{2}{l}{range} & \multicolumn{2}{l}{$0.0050$}\\
\midrule
\multicolumn{4}{l}{\emph{compressed distance priors}}\\
\citet{yang2025}, five datasets                    & priors, DESI DR1 & $-0.00612$ & $0.00243$\\
this work, configuration (a)                       & priors           & $+0.0060$ & $0.0030$\\
this work, configuration (b)                       & priors           & $-0.0030$ & $0.0010$\\
this work, configuration (c)                       & priors           & $-0.0100$ & $0.0020$\\
this work, direct $r_s$, $\mathcal{R}$ label       & priors           & $+0.0020$ & $0.0020$\\
\midrule
\multicolumn{2}{l}{range} & \multicolumn{2}{l}{$\mathbf{0.0160}$}\\
\bottomrule
\end{tabular}
\caption{$^\dagger$ Approximate conversion: this model's $Q$ is also sourced by baryons and
radiation. We include it because it \emph{widens} the full-likelihood family and therefore
makes our own point harder to make. $^\ddagger$ Nested predecessor of \citet{li2025}, same
group; included and flagged. Conversions: $\epsilon=3\nu f$ or $\epsilon=-3w_{\rm dm}f$ with
$f=\rho_{\rm dm}/\rho_m=0.8389$.}
\label{tab:litterature}
\end{table}

The full-likelihood family is internally consistent ($\chi^2/\mathrm{dof}=1.87$ for five
degrees of freedom) and spans $0.0050$; the compressed-prior family spans $0.0160$, a factor
$3.23$. Excluding the flagged approximate conversion --- the comfortable choice, which we do
not take --- the factor would be $4.67$ and $\chi^2/\mathrm{dof}=0.72$.

Three caveats we grant in full. The six full-likelihood models are not identical to
\eqref{eq:famille}: \citet{tsiapi2019} let $\rho_\Lambda$ run and carry a
$\Omega_{\rm dm}/(1-\nu)$ prefactor in $E^2$, so at equal $\epsilon$ their $H(z)$ is not ours;
\citet{kumar2025} and \citet{yaoliu2025} attach different perturbation prescriptions
($c_s^2=0$ versus $c_s^2=w$) to the same background, and their intervals are not transferable
to a background-only law. Their datasets differ from ours and from each other. And two of the
six are the same group on nested data. Table~\ref{tab:litterature} therefore measures a
\emph{dispersion}, and that is all we claim from it.

We also record, because it is the single most instructive item in the compressed column:
\citet{yang2025}'s headline $\epsilon=-0.0073^{+0.0029}_{-0.0033}$ at $2.4\sigma$ requires
five datasets; on DESI$+$CMB$+$CC alone the same paper reports
$\epsilon=+0.0023^{+0.0055}_{-0.0067}$ and concludes there is no significant deviation. A
sign reversal inside a single publication, on nested data.

\section{What we do not claim}
\label{sec:limites}

We do not measure $\epsilon$. None of our four configurations is a full likelihood, and the
one that drops $\mathcal{R}$ is driven by the $H_0$ prior. We do not refute
\citet{yang2025} or anyone else: our pipeline, data and nuisance treatment differ. We do not
claim that the field has abandoned $\mathcal{R}$ --- DESI DR2 conditions its baseline on
$(\theta_\ast,\omega_b,\omega_{bc})$ but explicitly tests the $(\mathcal{R},\ell_A,\omega_b)$
compression and reports it equivalent for the late-time models they consider
\citep{desidr2}. Dropping $\mathcal{R}$ is justified here narrowly, because $\sqrt{\Omega_m}$
has no definition when matter does not dilute as $a^{-3}$.

We also disclose a systematic of our own. The distance priors we use are those of
\citet{zhaiwang2019}, whose $\ell_A=301.80845$ corresponds to $100\theta=1.040923$, i.e.\ the
$\theta_{\rm MC}$ convention. Our $\chi^2$ predicts $\ell_A$ from CAMB's exact $z_\ast$ and
$r_\ast$, i.e.\ in the $\theta_\ast$ convention, yielding $\ell_A=301.75963$ and
$100\theta=1.041091$ at the $\Lambda$CDM optimum. The two conventions differ by
$-0.54\sigma$ in $\ell_A$ in every evaluation. This offset is common to all four
configurations and therefore does not generate the spread of Table~\ref{tab:quatre}, but it
is present and we state it rather than let a reader discover it.

\section{Recommendations}
\label{sec:reco}

\begin{enumerate}
\item A published constraint on a matter dilution exponent obtained from compressed CMB
      priors should state \emph{which} $\omega_m$ feeds $\mathcal{R}$, $r_d$, $r_\ast$ and
      $z_\ast$, and should quote the systematic obtained by making the other choice. In this
      family that systematic is of order $10^{-2}$ --- an order of magnitude above the
      statistical error usually quoted.
\item $\epsilon$ should be applied to the dark sector only, and the sound-horizon integrand
      should use the same baryon law as the background (\S\ref{sec:baryons}).
\item Where the compressed route is kept, $r_s$ should be integrated in the model's own
      background rather than taken from a $\Lambda$CDM-calibrated fitting formula. This is
      standard practice in the early-dark-energy literature and costs nothing; we validate an
      implementation to $1.3\times10^{-6}$ against CAMB in \S\ref{sec:mesure}.
\item Finally, we record a gap. A systematic search of the matter-creation literature
      returned \emph{no} analysis of any such model against full CMB power spectra: every one
      uses supernovae, BAO, $H(z)$, growth data, or a compressed CMB. The most recent example
      we found still describes its CMB input as ``a compressed (also referred to as
      geometrical or background) CMB likelihood, where the CMB is treated as a BAO
      measurement at $z\approx1100$'' \citep{schiavone2026}. Given \S\ref{sec:pouvoir}, the
      cost of that choice is not a loss of precision --- it is the introduction of a
      referent that the data cannot fix.
\end{enumerate}

\section*{Data and reproducibility}
Every number in this paper is produced by a script whose success criteria were frozen by
hash before execution, and re-verified afterwards by an independent tool that was itself
mutation-tested (five injected faults, five detected). Where a frozen criterion later proved
ill-posed, the criterion was left untouched and the defect recorded, rather than the
criterion rewritten.

\begin{thebibliography}{99}
\bibitem[Alcaniz \& Lima(2005)]{alcaniz2005} Alcaniz, J.~S., \& Lima, J.~A.~S. 2005,
  Phys.\ Rev.\ D, 72, 063516, \texttt{astro-ph/0507372}
\bibitem[DESI Collaboration(2025)]{desidr2} DESI Collaboration 2025, Phys.\ Rev.\ D, 112,
  083515, \texttt{arXiv:2503.14738}
\bibitem[Eisenstein \& Hu(1998)]{eh1998} Eisenstein, D.~J., \& Hu, W. 1998, ApJ, 496, 605
\bibitem[Elgar\o y \& Multam\"aki(2007)]{elgaroy2007} Elgar\o y, \O., \& Multam\"aki, T.
  2007, A\&A, 471, 65, \texttt{astro-ph/0702343}
\bibitem[Hu \& Sugiyama(1996)]{hs1996} Hu, W., \& Sugiyama, N. 1996, ApJ, 471, 542
\bibitem[Kumar et al.(2025)]{kumar2025} Kumar, U., Ajith, A., \& Verma, A. 2025,
  \texttt{arXiv:2504.14419}
\bibitem[Li et al.(2024)]{li2024} Li, T.-N., Wu, P.-J., Du, G.-H., et al. 2024, ApJ, 976, 1,
  \texttt{arXiv:2407.14934}
\bibitem[Li et al.(2025)]{li2025} Li, T.-N., et al. 2025, \texttt{arXiv:2510.11363}
\bibitem[Lima et al.(1996)]{lima1996} Lima, J.~A.~S., Germano, A.~S.~M., \& Abramo, L.~R.~W.
  1996, Phys.\ Rev.\ D, 53, 4287, \texttt{gr-qc/9511006}
\bibitem[Planck Collaboration(2020)]{planck2018} Planck Collaboration 2020, A\&A, 641, A6,
  \texttt{arXiv:1807.06209}
\bibitem[Schiavone et al.(2026)]{schiavone2026} Schiavone, T., De Angelis, M., Escamilla,
  L.~A., Montani, G., \& Di Valentino, E. 2026, \texttt{arXiv:2601.14222}
\bibitem[Tsiapi \& Basilakos(2019)]{tsiapi2019} Tsiapi, P., \& Basilakos, S. 2019, MNRAS,
  485, 2505, \texttt{arXiv:1810.12902}
\bibitem[Yang et al.(2025)]{yang2025} Yang, W., Dai, W.-B., \& Wang, D. 2025,
  \texttt{arXiv:2505.09879}
\bibitem[Yao \& Liu(2025)]{yaoliu2025} Yao, Y., \& Liu, X. 2025, \texttt{arXiv:2507.00478}
\bibitem[Zhai \& Wang(2019)]{zhaiwang2019} Zhai, Z., \& Wang, Y. 2019, JCAP, 07, 005,
  \texttt{arXiv:1811.07425}
\end{thebibliography}

\end{document}
